\documentclass{article}
\usepackage[utf8]{inputenc}
\usepackage[T1]{fontenc}
\usepackage{geometry}
\geometry{
    a4paper,
    left=2.5cm,
    right=2.5cm,
    top=2.5cm,
    bottom=2.5cm
}
\usepackage{graphicx}
\usepackage{float}
\usepackage{hyperref}
\usepackage[english]{babel}
\usepackage{amsmath}
\usepackage{amssymb}
\usepackage{array}
\usepackage{booktabs}
\usepackage{caption}
\usepackage{subcaption}
\usepackage{enumitem}
\usepackage{fancyhdr}
\usepackage{titlesec}
\usepackage{url}
\usepackage{textcomp}
\usepackage{xcolor}
\usepackage{listings}
\usepackage{multirow}
\usepackage{longtable}
\usepackage{tabularx}
\usepackage{times}

% Set up page style
\pagestyle{fancy}
\fancyhf{}
\rfoot{\thepage}
\renewcommand{\headrulewidth}{0pt}

% Title formatting
\titleformat{\chapter}[display]{\normalfont\bfseries}{\chaptertitle\thechapter}{0pt}{\Huge}
\titlespacing*{\chapter}{0pt}{50pt}{40pt}

% Document info
\title{Template Document}
\author{Academic Support Agent}
\date{\today}

\begin{document}
\maketitle

\maketitle

%  
  \LaTeX   Assignment    .205	 		
DSI205	Basics of Linear Algebra for Data Analytics

\section{  Latex   PC}
\begin{enumerate}
\item     Install \code{MiKTex}   \url{https://miktex.org/}   \code{TeXworks}   Editor \\
\code{MiKTex}   Compiler   Latex   PDF   \code{TeXworks}   Editor   Latex  
\item   \url{https://www.overleaf.com/read/hyhynwsnpjdx}  \\
  Menu > Download > Source
\item   \code{.tex}   \code{TeXworks}   Compiler   \textbf{XeLaTex}   Menu
\item     Compile/Typeset
\item   Install   Library
\item   PDF  
\end{enumerate}


\section{  Overleaf   Online}
  Latex     \emph{Overleaf} \url{https://www.overleaf.com} \\
 :   \emph{Overleaf}   Compile  
\begin{enumerate}
\item   \href{https://www.overleaf.com}{Overleaf}
\item   \url{https://www.overleaf.com/latex/templates/dsi205-template/tsnyrzxnngdw}   Open as Template (  Overleaf  )
\item   Menu > Compiler   Compiler   \textbf{XeLaTex}
\item   Recompile   Auto-Compile (      Auto Compile   On)
\item     Recompilie   PDF   Menu > Download > PDF
\end{enumerate}

\section{ }
  Math Mode   3  

\subsection{Inline mode}
    $\backslash( \dots \backslash)$, $\$ \dots \$$,   $\backslash$begin\{math\} $\dots \backslash$end\{math\}  

\code{\$E = mc\^2\$}   $E = mc^2$

\code{$\backslash($ E = mc\^2 $\backslash)$}   \( E = mc^2 \)

\subsection{Displayed Mode   Unnumbered}
      $\backslash[ \dots \backslash]$   \$\$ $\dots$ \$\$  

\code{$\backslash[$ $\backslash$bar\{x\} = $\backslash$sum x\_i $\backslash]$}  
\[ \bar{x} = \sum x_i \]

\code{\$\$ $\backslash$bar\{x\} = $\backslash$sum x\_i \$\$}  
$$ \bar{x} = \sum x_i $$

\subsection{Displayed Mode   Numbered}
          $\backslash$begin\{equation\} $\dots \backslash$end\{equation\}

  \\
$\backslash$begin\{equation\} \\
I = $\backslash$frac\{V\}\{R\} \\
$\backslash$end\{equation\}

 
\begin{equation}
I = \frac{V}{R}
\end{equation}

\section{ }
  \code{aligned}   \code{equation}
\begin{equation}
\begin{aligned}
a & = 2 \times 2 \times b \times c \\
a & = 4 \times b \times c \\
a & = 4bc
\end{aligned}
\end{equation}

\section{   }
  \code{amsmath}   \code{bmatrix}   \code{equation}
\begin{equation}
{\mathbf v} = \begin{bmatrix}1 \\ 5 \\ -4 \\ 0\end{bmatrix}
\end{equation}

\begin{equation}
A = \begin{pmatrix}1 & 2 & -3 \\ 4 & 0 & 1\end{pmatrix}
\end{equation}

\begin{equation}
B = \begin{bmatrix}1 & 2 & -3 \\ 4 & 0 & 1\end{bmatrix}
\end{equation}

\begin{equation}
c = \det A = \begin{vmatrix}-2 & 0 & 1 \\ 5 & -4 & 2 \\ 3 & 1 & 8\end{vmatrix}
\end{equation}
\end{document}